% !PRINTED none
% [part-title page: no running-head folio visible on this scan; opens the Third Section and § 27]
\begin{center}
\gesperrt{Third Section.}

\textbf{\gesperrt{Transformation of the Theta-Functions.}}
\end{center}

\noindent\hrulefill

\begin{center}
\textbf{§ 27. \gesperrt{The Principle of Transformation.}}
\end{center}

We return to the defining equations of the $T$-functions of
order $m$ given in § 17, and there, for a reason that will be evident
at once, we furnish the letters $\omega$, $a$, $b$, $m$
with accents, so that these equations read:
\weq{1}{\begin{gathered}
T(u+\omega'_1) = e^{-\pi i [a'_1(2u+\omega'_1)+b'_1]} T(u), \\
T(u+\omega'_2) = e^{-\pi i a'_2 (2u+\omega'_2)+b'_2} T(u). % [sic] German source prints the exponent without the expected brackets.
\end{gathered}}
\weq{2}{a'_2\omega'_1 - a'_1\omega'_2 = m'.}
If $a$, $c$ are any integers whatever (positive or negative),
then, on replacing $n_1$, $n_2$ in § 17, (16) by $-c$, $a$, there results:
\weq{3}{\begin{gathered}
T(u - c\omega'_1 + a\omega'_2) \\
= e^{-\pi i(-ca'_1+aa'_2)(2u-c\omega'_1+a\omega'_2) - \pi i(-cb'_1+ab'_2-m'ac)} T(u),
\end{gathered}}
an equation which also follows from one of the equations (1)
if in it $a'$, $b'$, $\omega'$ are replaced by
\[
-ca'_1+aa'_2,\quad -cb'_1+ab'_2-m'ac,\quad -c\omega'_1+a\omega'_2
\]
respectively.

Herein is contained the principle of the transformation of the
$T$-functions.

Let $b$, $\partial$ be two other integers for which the
determinant
\weq{4}{n = a\partial - bc}
has a \gesperrt{positive} value. We put
\weq{5}{\begin{gathered}
\omega_1 = +\partial\omega'_1 - b\omega'_2, \\
\omega_2 = -c\omega'_1 + a\omega'_2,
\end{gathered}}
and consequently
\weq{6}{\begin{gathered}
n\omega'_1 = a\omega_1 + b\omega_2 \\
n\omega'_2 = c\omega_1 + \partial\omega_2.
\end{gathered}}
